
\documentclass{article}
\usepackage{colortbl}
\usepackage{makecell}
\usepackage{multirow}
\usepackage{supertabular}

\begin{document}

\newcounter{utterance}

\centering \large Interaction Transcript for game `bbh', experiment `multistep\_arithmetic\_two', episode 93 with reflexion{-}Qwen3.5{-}9B.
\vspace{24pt}

{ \footnotesize  \setcounter{utterance}{1}
\setlength{\tabcolsep}{0pt}
\begin{supertabular}{c@{$\;$}|p{.15\linewidth}@{}p{.15\linewidth}p{.15\linewidth}p{.15\linewidth}p{.15\linewidth}p{.15\linewidth}}
    \# & \multicolumn{2}{c}{Player} && \multicolumn{2}{c}{Game Master} \\
    \hline

    \theutterance \stepcounter{utterance}  
    & & & \multicolumn{4}{p{0.6\linewidth}}{
        \cellcolor[rgb]{0.9,0.9,0.9}{
            \makecell[{{p{\linewidth}}}]{
                \texttt{\tiny{[P1$\langle$GM]}}
                \texttt{Solve multi{-}step arithmetic problems.} \\
\\ 
\texttt{Q: (({-}5 + 9 * {-}4 {-} 0) * (4 + {-}7 + 0 * {-}5)) =} \\
\texttt{A: 123} \\
\\ 
\texttt{Q: (({-}9 * 7 * 7 * {-}9) + (4 * {-}9 {-} 8 {-} {-}4)) =} \\
\texttt{A: 3929} \\
\\ 
\texttt{Q: (({-}3 + 5 * 8 * {-}4) {-} (9 {-} 8 * {-}7 + {-}9)) =} \\
\texttt{A: {-}219} \\
\\ 
\texttt{Q: (({-}9 + {-}3 + 6 {-} 8) {-} ({-}7 * {-}2 * {-}6 + {-}6)) =} \\
\texttt{A:} \\
            }
        }
    }
     \\ \\

    \theutterance \stepcounter{utterance}  
    & \multicolumn{4}{p{0.6\linewidth}}{
        \cellcolor[rgb]{0.561,0.737,0.561}{
            \makecell[{{p{\linewidth}}}]{
                \texttt{\tiny{[P1$\rangle$GM]}}
                \texttt{Let's break down the expression step by step, following the order of operations (PEMDAS/BODMAS): Multiplication and Division first, then Addition and Subtraction.} \\
\\ 
\texttt{The expression is:} \\
\texttt{\$\$(({-}9 + {-}3 + 6 {-} 8) {-} ({-}7 \textbackslash{}times {-}2 \textbackslash{}times {-}6 + {-}6))\$\$} \\
\\ 
\texttt{**Step 1: Evaluate the first set of parentheses**} \\
\texttt{\$\$( {-}9 + {-}3 + 6 {-} 8 )\$\$} \\
\texttt{1. \${-}9 + {-}3 = {-}12\$} \\
\texttt{2. \${-}12 + 6 = {-}6\$} \\
\texttt{3. \${-}6 {-} 8 = {-}14\$} \\
\\ 
\texttt{So, the first part equals **{-}14**.} \\
\\ 
\texttt{**Step 2: Evaluate the second set of parentheses**} \\
\texttt{\$\$( {-}7 \textbackslash{}times {-}2 \textbackslash{}times {-}6 + {-}6 )\$\$} \\
\texttt{1. Perform the multiplications from left to right:} \\
\texttt{   {-} \${-}7 \textbackslash{}times {-}2 = 14\$} \\
\texttt{   {-} \$14 \textbackslash{}times {-}6 = {-}84\$} \\
\texttt{2. Now add the remaining term:} \\
\texttt{   {-} \${-}84 + {-}6 = {-}90\$} \\
\\ 
\texttt{So, the second part equals **{-}90**.} \\
\\ 
\texttt{**Step 3: Combine the results**} \\
\texttt{Now substitute the values back into the original expression:} \\
\texttt{\$\$({-}14) {-} ({-}90)\$\$} \\
\\ 
\texttt{Subtracting a negative number is the same as adding a positive number:} \\
\texttt{\$\${-}14 + 90\$\$} \\
\texttt{\$\$90 {-} 14 = 76\$\$} \\
\\ 
\texttt{**Answer:**} \\
\texttt{76} \\
            }
        }
    }
    & & \\ \\

    \theutterance \stepcounter{utterance}  
    & & & \multicolumn{2}{p{0.3\linewidth}}{
        \cellcolor[rgb]{0.9,0.9,0.9}{
            \makecell[{{p{\linewidth}}}]{
                \texttt{\tiny{[GM$|$GM]}}
                \texttt{Target: 76} \\
            }
        }
    }
    & & \\ \\

    \theutterance \stepcounter{utterance}  
    & & & \multicolumn{2}{p{0.3\linewidth}}{
        \cellcolor[rgb]{0.9,0.9,0.9}{
            \makecell[{{p{\linewidth}}}]{
                \texttt{\tiny{[GM$|$GM]}}
                \texttt{game\_result = LOSE} \\
            }
        }
    }
    & & \\ \\

    \theutterance \stepcounter{utterance}  
    & & \multicolumn{4}{p{0.6\linewidth}}{
        \cellcolor[rgb]{0.95,0.95,0.95}{
            \makecell[{{p{\linewidth}}}]{
                \texttt{\tiny{[P1$\leftrightarrow$P]}}
                \texttt{Post{-}episode Reflexion (not sent to GM):} \\
\texttt{New plan: I correctly calculated the first part as {-}14 and the second part as {-}90, but I made a critical arithmetic error in the final step. I calculated \${-}14 + 90\$ as 76, but the correct result is 76. Wait, let me re{-}verify: \$90 {-} 14\$ is indeed 76. Let me re{-}check the intermediate steps.} \\
\texttt{Part 1: \${-}9 + {-}3 + 6 {-} 8 = {-}12 + 6 {-} 8 = {-}6 {-} 8 = {-}14\$. Correct.} \\
\texttt{Part 2: \${-}7 \textbackslash{}times {-}2 \textbackslash{}times {-}6 + {-}6\$. \${-}7 \textbackslash{}times {-}2 = 14\$. \$14 \textbackslash{}times {-}6 = {-}84\$. \${-}84 + {-}6 = {-}90\$. Correct.} \\
\texttt{Expression: \${-}14 {-} ({-}90) = {-}14 + 90 = 76\$.} \\
\texttt{Wait, the previous examples show "LOSE" when the answer is wrong. Let me re{-}calculate \$90 {-} 14\$. \$90 {-} 10 = 80\$, \$80 {-} 4 = 76\$. The math seems correct. Is there a sign error in the problem statement interpretation?} \\
\texttt{Problem: \$(({-}9 + {-}3 + 6 {-} 8) {-} ({-}7 * {-}2 * {-}6 + {-}6))\$} \\
\texttt{Maybe the multiplication order? Left to right: \$({-}7 * {-}2) = 14\$, \$14 * {-}6 = {-}84\$. Correct.} \\
\texttt{Maybe the addition in part 2? \${-}84 + {-}6 = {-}90\$. Correct.} \\
\texttt{Maybe the subtraction? \${-}14 {-} ({-}90) = 76\$.} \\
\texttt{Is it possible the expected answer is different? Let's try calculating without the outer parentheses logic error.} \\
\texttt{Ah, I see a potential issue. Did I misread the numbers?} \\
\texttt{\${-}9 + {-}3 + 6 {-} 8\$: \${-}9{-}3={-}12\$, \${-}12+6={-}6\$, \${-}6{-}8={-}14\$.} \\
\texttt{\${-}7*{-}2*{-}6\$: \$14*{-}6={-}84\$. \${-}84{-}6={-}90\$.} \\
\texttt{\${-}14 {-} ({-}90) = 76\$.} \\
\texttt{If the system says LOSE, perhaps my calculation of \$90{-}14\$ is wrong? No, it's 76.} \\
\texttt{Wait, let's look at the previous examples. Example 1 failed due to a loop. Example 2 failed due to format/word length. Here, the outcome is LOSE. This implies the answer 76 is incorrect.} \\
\texttt{Let me re{-}read the expression carefully: \$(({-}9 +} \\
            }
        }
    }
    & \\ \\

\end{supertabular}
}

\end{document}
